\documentclass[12pt]{article}
\usepackage[T1]{fontenc}
\usepackage[utf8]{inputenc}
\usepackage{lmodern}
\usepackage{textcomp}
\usepackage{enumitem}
\usepackage{geometry}
\usepackage{graphicx}
\usepackage{amsmath, amssymb}
\geometry{margin=1in}
\begin{document}
\begin{center}
History - Graduate Level Assessment 11 \
Total Marks: 40 \
Answer all questions.
\end{center}

\section*{Multiple Choice (10 marks)}
\begin{enumerate}[label=\arabic*.]
\item Research on C 3 and C 4 pathways indicates that changes in plant communities were associated with a transition in the diets of ancient hominids from:
\begin{enumerate}[label=(\alph*)]
    \item meat to fruits and vegetables.
    \item species focused on grasses to fruits and vegetables.
    \item nuts and fruits to fish.
    \item nuts and fruits to species more focused on grasses.
    \item species focused on grasses to meat.
\end{enumerate}
\item Markings found on the mud-bricks of Huaca del Sol are probably:
\begin{enumerate}[label=(\alph*)]
    \item instructions for the builders written there by the engineers.
    \item random scratches and not a meaningful script or written language.
    \item intended to identify the local group that had produced them.
    \item simple arithmetic performed by carpenters on the building site.
\end{enumerate}
\item Archaeological evidence indicates that cattle were first domesticated where and how long ago?
\begin{enumerate}[label=(\alph*)]
    \item in sub-Saharan Africa, about 10,500 years ago
    \item in western Europe, about 3,500 years ago
    \item in the Middle East, about 10,500 years ago
    \item in North America, about 5,500 years ago
    \item in East Asia, about 5,500 years ago
\end{enumerate}
\item Unlike most other early civilizations, Minoan culture shows little evidence of:
\begin{enumerate}[label=(\alph*)]
    \item trade.
    \item warfare.
    \item the development of a common religion.
    \item conspicuous consumption by elites.
\end{enumerate}
\item Homo erectus differed from Homo habilis in which way?
\begin{enumerate}[label=(\alph*)]
    \item Erectus was not bipedal.
    \item Erectus was not capable of using tools.
    \item Erectus was primarily herbivorous.
    \item Erectus fossils are found only in Africa.
    \item Erectus possessed a larger brain.
\end{enumerate}
\end{enumerate}

\section*{True / False (10 marks)}
\textit{Answer True or False}
\begin{enumerate}[label=\arabic*.]
\setcounter{enumi}{5}
\item True or False: In 2005, the company sold its personal computer business to Chinese technology company Lenovo, and in the same year it agreed to acquire Micromuse. A year later IBM launched Secure Blue, a low-cost hardware design for data encryption that can be built into a microprocessor. In 2009 it acquired software company SPSS Inc.
\item True or False: Warsaw is located on two main geomorphologic formations: the plain moraine plateau and the Vistula Valley with its asymmetrical pattern of different terraces. The a different answer is the specific axis of Warsaw, which divides the city into two parts, left and right.
\item True or False: There was a special case established under the State Management Scheme where the brewery and licensed premises were bought and run by the state until 1973, most notably in a different answer.
\item True or False: Rome banned it on several occasions under extreme penalty. A law passed in 81 BC characterised human sacrifice as a different answer committed for magical purposes.
\item True or False: In the latter part of the second revolution, Thomas Alva Edison developed many devices that greatly influenced life around the world and is often credited with the creation of the first industrial research laboratory. In 1882, Edison switched on a different answer that provided 110 volts direct current to fifty-nine customers in lower Manhattan.
\end{enumerate}

\section*{Long Form Response (20 marks)}
\textit{Answer each question with a clear and complete explanation.}
\begin{enumerate}[label=\arabic*.]
\setcounter{enumi}{10}
\item Read the following passage and answer the question: The county was established in 1182, later than many other counties. During Roman times the area was part of the Brigantes tribal area in the military zone of Roman Britain. The towns of Manchester, Lancaster, Ribchester, Burrow, Elslack and Castleshaw grew around Roman forts. In the centuries after the Roman withdrawal in 410 AD the northern parts of the county probably formed part of the Brythonic kingdom of Rheged, a successor entity to the Brigantes tribe. During the mid-8 th century, the area was incorporated into the Anglo-Saxon Kingdom of Northumbria, which became a part of England in the 10 th century. Question: When was the area part of Brigantes tribal area?
\item Read the following passage and answer the question: The very large and ornate School Hall and School Library (by L K. Hall) were erected in 1906-8 across the road from Upper School as the school's memorial to the Etonians who had died in the Boer War. Many tablets in the cloisters and chapel commemorate the large number of dead Etonians of the Great War. A bomb destroyed part of Upper School in World War Two and blew out many windows in the Chapel. The college commissioned replacements by Evie Hone (1949-52) and by John Piper and Patrick Reyntiens (1959 onwards). Question: Who was hired to replace windows in the Chapel from 1949-1952?
\end{enumerate}

\vfill
\noindent\textit{End of Paper}
\end{document}